\documentclass[11pt,letterpaper]{article}
\usepackage[margin=0.72in]{geometry}
\usepackage[T1]{fontenc}
\usepackage{lmodern}
\usepackage{amsmath,amssymb}
\usepackage{enumitem}
\usepackage[hidelinks]{hyperref}
\hypersetup{pdftitle={Practice Questions for Midterm 2},pdfauthor={Manish Patnaik}}
\setlength{\parindent}{0pt}
\setlength{\parskip}{3pt}
\setlist[enumerate]{leftmargin=*,topsep=2pt,itemsep=2pt,parsep=0pt}
\newcommand{\question}[1]{\par\medskip\noindent\textbf{#1}\par}

\begin{document}
\begin{center}
{\Large\bfseries Practice Questions for Midterm 2}\\[2pt]
MATH 327: Algebra 1 \quad\textbar\quad Fall 2026
\end{center}
\vspace{-5pt}
\hrule
\vspace{5pt}
\textbf{Directions.} Show your work and justify your claims. State any theorem you use and check its hypotheses. You may use results proved in class through Lecture 23 unless a question asks you to prove that result.

% Adapted from 2025-Fall-MATH327-MidtermExam.tex, Question 3.
\question{1. Products and quotients}
Let $H,K\leq G$, and define $HK=\{hk:h\in H,\ k\in K\}$. Assume that $H$ normalizes $K$: $hKh^{-1}=K$ for every $h\in H$.
\begin{enumerate}[label=(\alph*)]
\item Show that $HK$ is a subgroup of $G$.
\item Show that $K\trianglelefteq HK$ and $H\cap K\trianglelefteq H$.
\item Construct an isomorphism $H/(H\cap K)\cong HK/K$. Prove that your map is well defined and is an isomorphism.
\end{enumerate}

% Based on the S_4 -> S_3 action and correspondence theorem in Lectures 15--18.
\question{2. The homomorphism $S_4\to S_3$}
Let $S_4$ act on the three partitions $\Pi_1=12\mid34$, $\Pi_2=13\mid24$, and $\Pi_3=14\mid23$, where $12\mid34$ abbreviates $\{\{1,2\},\{3,4\}\}$. Let $\varphi:S_4\to S_3$ be the resulting homomorphism. You may use that $\varphi$ is surjective and that $\ker\varphi=V=\{1,(12)(34),(13)(24),(14)(23)\}$. Set $H=\varphi^{-1}(\langle(13)\rangle)$.
\begin{enumerate}[label=(\alph*)]
\item Compute $\varphi((1234))$.
\item Show that $H=\langle(1234),(12)(34)\rangle$, and determine $|H|$.
\item Identify $H$ up to isomorphism. Is $H$ normal in $S_4$? Justify both answers.
\end{enumerate}

% Adapted from 2025-Fall-MATH327-MidtermExam.tex, Question 4.
\question{3. Conjugation in $S_3$}
Let $G$ be a group. Define $g\cdot x=gxg^{-1}$ for $g,x\in G$.
\begin{enumerate}[label=(\alph*)]
\item Show that this rule defines a left action of $G$ on itself.
\item Show that the center $Z(G)$ is the union of all orbits of cardinality $1$.
\item Show that every normal subgroup of $G$ is a union of orbits under this action.
\item Describe all orbits and stabilizers for this action when $G=S_3$.
\end{enumerate}

% Adapted from 2026-Summer-MATH327-MidtermExam2-corrected.tex, Question 1.
\question{4. Conjugation in $\mathrm{GL}_2(\mathbb R)$}
Let $A=\begin{pmatrix}1&0\\0&2\end{pmatrix}$. The group $\mathrm{GL}_2(\mathbb R)$ acts on itself by conjugation: $g\cdot B=gBg^{-1}$.
\begin{enumerate}[label=(\alph*)]
\item Describe the stabilizer of $A$ under this action. Justify your answer.
\item Describe the orbit of $A$ under this action. Justify your answer.
\end{enumerate}

% Adapted from 2026-Summer-MATH327-MidtermExam2-corrected.tex, Question 2.
\question{5. A normal subgroup of order $5$}
Let $N\trianglelefteq G$, where $|N|=5$ and $|G|$ is odd. Prove that $N\subseteq Z(G)$.
\end{document}
